Background section:

We begin with explaining that the following discussion is entirely focused on single-node shared-memory systems and does not deal with/extrapolate fully to distributed systems.

1. Data Layout:
	1.1 In-memory:
		(a) Adj Matrix
		(b) AdjList
		(c) CSR
		(d) edgelists <-- why?

	1.2 Persistent Layouts:
		(a) Transactional Edge Log (weakest form of persistence: MMAP based)
		(b) Neo4J : on-disk layout + linked lists in memory
		(c)	Preprocessed files + sharded files: Blaze, Bolt, GraphChi, X-Stream
		(d) Proper DB: fit stuff into a Tree-based layouts: LSMT, B-Trees

2. Workloads:
	2.1 OLTP-like:
		(a) Characteristics: R/W mix, insert/delete/short scans
		(b) Requirements: efficient indexing to find neighborhoods, sequential scans, good join algorithms
		(c) Special storage requirements for the properties of graph elements:
			(i) Labeled Property Graph
		(d) LDBC is used for benchmarking

	2.2 OLAP-like:
		(a) Characteristics: Read-only, snapshot-based, needs high bandwidth for whole graph traversals
		(b) Kernels: The different types of benchmark kernels and why they are interesting. 
		(c) Need access to whole graph and the computed value is defined for each vertex (vertex functions)

3. Computational Model:
	- Graph processing algorithms can be defined as matrix multiplications: Y = X'A (X' is transpose vector matrix). Under this framing, we need to access X' and load in the relevant "rows" in A (out or in-neighborhood) to produce the Y vector (value vector of cardinality |V|). X' can be placed in memory, since |V| is of a manageable size.
	
	- In a single-node system, the computation model defines:
		(i) [IO]is A read in its entirety each time, or is there some reasonable way to access a subgraph, ideally only the neighborhood needed (all: naive linear algebra, whole matrix in a smart way: x-stream, only relevant edges: most systems)

		(ii) [Consistency Semantics] Is Y immediately visible? There can be a synchronization barrier at the end of a round of computation and only after that are the results of a round available for use in the next round(Most systems). Otherwise, the values can be immediately applied and no synchronization is needed (GraphChi) // What does this mean for a single-node system?

		(iii) ?

	- Computational model dictates the data layout; the workload influences which data layout is useful and which computational model makes sense.

	- Examine the computational model used by the systems in 1.2(c) 